\documentclass[11pt,a4paper]{article}

% ============================================================
%  QIMMA -- COURS : Géométrie analytique de l'espace
%  (droites, plans, sphère)
%  2ème Bac Sciences Mathématiques (A et B)
%  Standard exhaustif : définitions formelles + démonstrations
%  Basé sur templates/qimma-template.tex (charte Qimma)
% ============================================================

\usepackage[french]{babel}
\usepackage[utf8]{inputenc}
\usepackage[T1]{fontenc}
\usepackage{lmodern}
\usepackage[a4paper,margin=2cm,top=2.3cm,bottom=2.6cm]{geometry}
\usepackage{amsmath,amssymb}
\usepackage{xcolor}
\usepackage{graphicx}
\usepackage{tikz}
\usetikzlibrary{shadows,calc}
\usepackage[most]{tcolorbox}
\usepackage{titlesec}
\usepackage{fancyhdr}
\usepackage{enumitem}
\usepackage{pifont}
\usepackage{array}
\usepackage{colortbl}
\usepackage{hyperref}

% ------------------- CHARTE QIMMA -------------------
\definecolor{qteal}{HTML}{0d9488}
\definecolor{qtealdark}{HTML}{134e4a}
\definecolor{qamber}{HTML}{fbbf24}
\definecolor{qambertext}{HTML}{92400e}
\definecolor{ink}{HTML}{1f2937}
\definecolor{inkgray}{HTML}{6b7280}
\definecolor{tealpale}{HTML}{e6f5f3}
\colorlet{colDef}{qteal}
\definecolor{colProp}{HTML}{0f766e}
\definecolor{colEx}{HTML}{b45309}
\definecolor{colRem}{HTML}{9d174d}
\color{ink}

\graphicspath{{../../../../assets/}}
\newcommand{\logofile}{../../../../assets/qimma-logo.pdf}

% ------------------- METADONNEES -------------------
\newcommand{\matiere}{Mathématiques}
\newcommand{\niveau}{2\up{ème} Bac -- Sciences Mathématiques}
\newcommand{\chapitretitre}{Géométrie analytique de l'espace}

% ------------------- EN-TETE / PIED DE PAGE -------------------
\pagestyle{fancy}
\fancyhf{}
\renewcommand{\headrulewidth}{0pt}
\renewcommand{\footrulewidth}{0.5pt}
\fancyfoot[L]{\raisebox{-0.15cm}{\includegraphics[height=0.35cm]{\logofile}}~\small\sffamily\color{inkgray}Qimma}
\fancyfoot[C]{\small\sffamily\color{inkgray}\matiere\ -- \niveau}
\fancyfoot[R]{\small\sffamily\color{qtealdark}\thepage}

% ------------------- TITRES DE SECTION -------------------
\titleformat{\section}
  {\normalfont\sffamily\Large\bfseries\color{qtealdark}}
  {\tikz[baseline=-3.2pt]{\node[circle,fill=qteal,text=white,inner sep=0pt,minimum size=8mm]{\sffamily\bfseries\thesection};}}
  {10pt}{}
  [\vspace{-2mm}{\color{qamber}\titlerule[1.6pt]}\vspace{2mm}]
\titlespacing*{\section}{0pt}{16pt}{10pt}

\titleformat{\subsection}
  {\normalfont\sffamily\large\bfseries\color{qtealdark}}
  {\color{qteal}\ding{212}}{6pt}{}
\titlespacing*{\subsection}{0pt}{12pt}{6pt}

% ------------------- BOITES PEDAGOGIQUES -------------------
\newtcolorbox{fiche}[3][]{
  enhanced, breakable,
  colback=white, colframe=#2,
  boxrule=1pt, arc=2mm,
  top=7mm, left=4mm, right=4mm, bottom=3mm,
  drop shadow={black!25, shadow xshift=1mm, shadow yshift=-1mm},
  attach boxed title to top left={xshift=4mm, yshift=-3mm, yshifttext=-1mm},
  boxed title style={colback=#2, arc=1.5mm, boxrule=0pt, left=2mm, right=2mm},
  coltitle=white, fonttitle=\sffamily\bfseries,
  title={#3}, #1
}
\newenvironment{definition}[1][Définition]
  {\begin{fiche}{colDef}{\ding{110}~~#1}}{\end{fiche}}
\newenvironment{propriete}[1][Propriété]
  {\begin{fiche}{colProp}{\ding{72}~~#1}}{\end{fiche}}
\newenvironment{exemple}[1][Exemple]
  {\begin{fiche}{colEx}{\ding{217}~~#1}}{\end{fiche}}
\newenvironment{remarque}[1][Remarque]
  {\begin{fiche}{colRem}{\ding{43}~~#1}}{\end{fiche}}

% Démonstration (hors boîte, style discret)
\newenvironment{demo}
  {\par\smallskip\noindent{\sffamily\bfseries\color{qtealdark}Démonstration.}\ \itshape}
  {\ \normalfont\hfill$\blacksquare$\par\medskip}

\hypersetup{colorlinks=true, linkcolor=qtealdark, urlcolor=qtealdark}

% ============================================================
\begin{document}

% ------------------- BANNIERE DE TITRE -------------------
\begin{tcolorbox}[
  enhanced, boxrule=0pt, arc=4mm,
  interior style={left color=qtealdark, right color=qteal},
  top=8mm, bottom=8mm, left=10mm, right=32mm,
  drop shadow={black!35, shadow xshift=1.5mm, shadow yshift=-1.5mm},
  before skip=0pt, after skip=9mm,
  overlay={
    \node[anchor=north west] at ([xshift=6mm,yshift=-6mm]frame.north west)
      {\includegraphics[height=1.25cm]{\logofile}};
    \node[anchor=north east, fill=qamber, text=qambertext, rounded corners=2pt,
          inner sep=4pt, font=\sffamily\bfseries\small] at ([xshift=-6mm,yshift=-6mm]frame.north east) {COURS};
  }
]
\hspace{1.6cm}\centering
{\sffamily\bfseries\color{white}\fontsize{23}{27}\selectfont \chapitretitre}\\[3mm]
{\sffamily\color{white!90}\large \matiere\ \textbullet\ \niveau}
\end{tcolorbox}

% ------------------- OBJECTIFS -------------------
\begin{tcolorbox}[
  enhanced, colback=tealpale, colframe=qteal, boxrule=0.8pt, arc=2mm,
  left=5mm, right=5mm, top=3mm, bottom=3mm,
  title={\sffamily\bfseries\color{qtealdark}\ding{51}~~Objectifs du chapitre},
  coltitle=qtealdark, colbacktitle=tealpale, boxrule=0.8pt,
  attach title to upper={\par\vspace{1mm}}
]
\begin{itemize}[leftmargin=6mm, label=\ding{212}, itemsep=1mm]
  \item Déterminer représentation paramétrique et équation cartésienne d'une droite ou d'un plan de l'espace.
  \item Étudier toutes les positions relatives : droite-droite, droite-plan, plan-plan.
  \item Calculer des distances (point-droite, point-plan) et déterminer des projections orthogonales.
  \item Calculer un angle entre deux droites, deux plans, ou une droite et un plan.
  \item Déterminer l'équation d'une sphère et étudier ses intersections avec une droite ou un plan.
  \item Résoudre les exercices type examen national combinant vecteurs, produit scalaire/vectoriel et sphères.
\end{itemize}
\end{tcolorbox}

\vspace{4mm}

% ------------------- SECTION 1 -------------------
\section{Représentations d'une droite}

\begin{propriete}[Représentation paramétrique]
La droite $\mathcal{D}$ passant par $A(x_0,y_0,z_0)$ de vecteur directeur $\vec u(a,b,c)\ne\vec0$ est l'ensemble des points $M(x,y,z)$ tels que
\[
  \begin{cases} x=x_0+ta \\ y=y_0+tb \\ z=z_0+tc \end{cases} \qquad t\in\mathbb{R}.
\]
\end{propriete}

\begin{demo}
$M\in\mathcal{D} \iff \overrightarrow{AM}$ colinéaire à $\vec u \iff \exists t\in\mathbb{R},\ \overrightarrow{AM}=t\vec u$, ce qui donne exactement le système ci-dessus en coordonnées.
\end{demo}

\begin{exemple}[Paramétrer une droite]
Droite passant par $A(1,-1,2)$, de vecteur directeur $\vec u(2,0,-1)$ :
\[
  \begin{cases} x=1+2t \\ y=-1 \\ z=2-t \end{cases} \qquad t\in\mathbb{R}.
\]
\end{exemple}

% ------------------- SECTION 2 -------------------
\section{Positions relatives de deux droites}

\begin{propriete}[Classification]
Deux droites $\mathcal{D}(A,\vec u)$ et $\mathcal{D}'(B,\vec v)$ de l'espace sont :
\begin{itemize}[leftmargin=6mm, itemsep=0.5mm]
  \item \textbf{parallèles} si $\vec u,\vec v$ colinéaires : \textbf{confondues} si de plus $A\in\mathcal{D}'$, sinon \textbf{strictement parallèles} ;
  \item \textbf{sécantes} si $\vec u,\vec v$ non colinéaires et $\overrightarrow{AB},\vec u,\vec v$ \textbf{coplanaires} (produit mixte nul) ;
  \item \textbf{non coplanaires} (gauches) si $\vec u,\vec v$ non colinéaires et $\overrightarrow{AB},\vec u,\vec v$ non coplanaires.
\end{itemize}
\end{propriete}

\begin{exemple}[Déterminer la position relative]
Soient $\mathcal{D}$ passant par $A(1,0,0)$, $\vec u(1,1,0)$, et $\mathcal{D}'$ passant par $B(0,1,1)$, $\vec v(1,0,1)$.

\smallskip
$\vec u,\vec v$ non colinéaires (coordonnées non proportionnelles). $\overrightarrow{AB}(-1,1,1)$.
\[
  \vec u\wedge\vec v = (1\times1-0\times0\,;\ 0\times1-1\times1\,;\ 1\times0-1\times1) = (1,-1,-1).
\]
\[
  \overrightarrow{AB}\cdot(\vec u\wedge\vec v) = (-1)(1)+(1)(-1)+(1)(-1) = -1-1-1 = -3 \ne 0.
\]
\textbf{Conclusion} : les droites sont \textbf{non coplanaires}.
\end{exemple}

% ------------------- SECTION 3 -------------------
\section{Plans : équation cartésienne, positions relatives}

\begin{remarque}[Rappel du chapitre précédent]
L'équation cartésienne d'un plan de vecteur normal $\vec n(a,b,c)$ passant par $A(x_0,y_0,z_0)$ est $a(x-x_0)+b(y-y_0)+c(z-z_0)=0$, soit $ax+by+cz+d=0$.
\end{remarque}

\begin{propriete}[Positions relatives de deux plans]
Soient $\mathcal{P}:ax+by+cz+d=0$ (normal $\vec n(a,b,c)$) et $\mathcal{P}':a'x+b'y+c'z+d'=0$ (normal $\vec n'(a',b',c')$).
\begin{itemize}[leftmargin=6mm, itemsep=0.5mm]
  \item $\mathcal{P}\parallel\mathcal{P}'$ (ou confondus) $\iff \vec n,\vec n'$ colinéaires.
  \item Sinon, $\mathcal{P}$ et $\mathcal{P}'$ sont \textbf{sécants} selon une droite, de vecteur directeur $\vec n\wedge\vec n'$.
\end{itemize}
\end{propriete}

\begin{propriete}[Position relative droite-plan]
Soit $\mathcal{D}(A,\vec u)$ et $\mathcal{P}$ de normal $\vec n$.
\begin{itemize}[leftmargin=6mm, itemsep=0.5mm]
  \item $\mathcal{D}\parallel\mathcal{P} \iff \vec u\cdot\vec n=0$ : \textbf{incluse} dans $\mathcal{P}$ si de plus $A\in\mathcal{P}$, sinon strictement parallèle.
  \item Sinon ($\vec u\cdot\vec n\ne0$), $\mathcal{D}$ est \textbf{sécante} à $\mathcal{P}$ en un unique point.
\end{itemize}
\end{propriete}

\begin{exemple}[Intersection droite-plan]
Soit $\mathcal{D}$ : $\begin{cases}x=1+t\\y=2-t\\z=t\end{cases}$ et $\mathcal{P}:x+y+z-4=0$. Déterminons leur intersection.

\smallskip
Substituons : $(1+t)+(2-t)+t-4=0\iff t-1=0\iff t=1$.

\smallskip
\textbf{Point d'intersection} : $x=2$, $y=1$, $z=1$, soit $M(2,1,1)$.
\end{exemple}

\begin{propriete}[Système de deux plans définissant une droite]
Deux plans \textbf{sécants} (non parallèles) $\mathcal{P}:ax+by+cz+d=0$ et $\mathcal{P}':a'x+b'y+c'z+d'=0$ définissent une droite $\mathcal{D}=\mathcal{P}\cap\mathcal{P}'$, de vecteur directeur $\vec n\wedge\vec n'$ (où $\vec n,\vec n'$ sont les normales respectives).
\end{propriete}

\begin{exemple}[Trouver la droite intersection de deux plans]
Soient $\mathcal{P}:x+y-z=0$ et $\mathcal{P}':2x-y+z-3=0$. Déterminons une représentation paramétrique de $\mathcal{P}\cap\mathcal{P}'$.

\smallskip
$\vec n(1,1,-1)$, $\vec n'(2,-1,1)$, non colinéaires : les plans sont sécants.
\[
  \vec n\wedge\vec n' = (1\times1-(-1)\times(-1)\,;\ (-1)\times2-1\times1\,;\ 1\times(-1)-1\times2) = (0,-3,-3).
\]
On peut simplifier le vecteur directeur en $(0,1,1)$.

\smallskip
\textbf{Point particulier} : cherchons un point du système en fixant $z=0$ : $x+y=0$ et $2x-y=3$. En additionnant : $3x=3\iff x=1$, puis $y=-1$. Point $A(1,-1,0)$.

\smallskip
\textbf{Conclusion} : $\mathcal{D}:\begin{cases}x=1\\y=-1+t\\z=t\end{cases}$, $t\in\mathbb{R}$.
\end{exemple}

\subsection{Angles entre droites et plans}

\begin{propriete}[Formules d'angle via le produit scalaire]
On définit l'angle géométrique (dans $\left[0\,;\frac\pi2\right]$) entre deux droites, deux plans, ou une droite et un plan, à l'aide de leurs vecteurs directeurs/normaux, en prenant la valeur absolue du cosinus (pour obtenir l'angle \textbf{aigu}, indépendant du sens choisi pour les vecteurs) :
\begin{itemize}[leftmargin=6mm, itemsep=0.5mm]
  \item \textbf{Angle entre deux droites} $\mathcal{D}(\vec u)$, $\mathcal{D}'(\vec u')$ :
  \[
    \cos\theta = \frac{|\vec u\cdot\vec u'|}{\|\vec u\|\,\|\vec u'\|}.
  \]
  \item \textbf{Angle entre deux plans} $\mathcal{P}(\vec n)$, $\mathcal{P}'(\vec n')$ :
  \[
    \cos\theta = \frac{|\vec n\cdot\vec n'|}{\|\vec n\|\,\|\vec n'\|}.
  \]
  \item \textbf{Angle entre une droite $\mathcal{D}(\vec u)$ et un plan $\mathcal{P}(\vec n)$} :
  \[
    \sin\theta = \frac{|\vec u\cdot\vec n|}{\|\vec u\|\,\|\vec n\|}.
  \]
\end{itemize}
\end{propriete}

\begin{remarque}[Pourquoi un sinus pour droite-plan]
L'angle entre une droite et un plan se mesure entre la droite et sa \textbf{projection} sur le plan, et non entre la droite et la normale : c'est pourquoi la formule utilise $\sin$ (angle complémentaire de celui entre $\vec u$ et $\vec n$) plutôt que $\cos$.
\end{remarque}

\begin{exemple}[Calcul d'angle entre deux plans]
Soient $\mathcal{P}:x+y+z-1=0$ et $\mathcal{P}':x-y+2=0$. Calculons l'angle entre ces deux plans.

\smallskip
$\vec n(1,1,1)$, $\vec n'(1,-1,0)$.
\[
  \cos\theta = \frac{|1\times1+1\times(-1)+1\times0|}{\sqrt3\times\sqrt2} = \frac{|0|}{\sqrt6} = 0.
\]
\textbf{Conclusion} : $\theta=\dfrac\pi2$ : les deux plans sont \textbf{perpendiculaires}.
\end{exemple}

% ------------------- SECTION 4 -------------------
\section{Distances et projections}

\begin{propriete}[Distance d'un point à une droite]
La distance du point $M_0$ à la droite $\mathcal{D}(A,\vec u)$ est
\[
  d(M_0,\mathcal{D}) = \frac{\left\|\overrightarrow{AM_0}\wedge\vec u\right\|}{\|\vec u\|}.
\]
\end{propriete}

\begin{demo}
$\left\|\overrightarrow{AM_0}\wedge\vec u\right\|$ est l'aire du parallélogramme construit sur $\overrightarrow{AM_0}$ et $\vec u$, qui vaut aussi (base $\times$ hauteur) $\|\vec u\|\times d(M_0,\mathcal{D})$ (la hauteur étant la distance cherchée). En isolant $d(M_0,\mathcal{D})$, on obtient la formule.
\end{demo}

\begin{exemple}[Calcul de distance point-droite]
Soit $\mathcal{D}$ passant par $A(0,0,0)$, $\vec u(1,1,1)$, et $M_0(1,0,0)$.
\[
  \overrightarrow{AM_0}(1,0,0), \qquad \overrightarrow{AM_0}\wedge\vec u = (0\times1-0\times1\,;\ 0\times1-1\times1\,;\ 1\times1-0\times1) = (0,-1,1).
\]
\[
  d(M_0,\mathcal{D}) = \frac{\sqrt{0+1+1}}{\sqrt3} = \frac{\sqrt2}{\sqrt3} = \frac{\sqrt6}{3}.
\]
\end{exemple}

% ------------------- SECTION 5 -------------------
\section{Sphères}

\begin{definition}[Sphère]
La \textbf{sphère} de centre $\Omega(x_0,y_0,z_0)$ et de rayon $R>0$ est l'ensemble des points $M(x,y,z)$ tels que $\Omega M=R$, d'équation
\[
  (x-x_0)^2+(y-y_0)^2+(z-z_0)^2 = R^2.
\]
\end{definition}

\begin{propriete}[Reconnaître une équation de sphère]
Une équation $x^2+y^2+z^2+ax+by+cz+d=0$ est celle d'une sphère si et seulement si, après mise sous forme canonique,
\[
  \left(x+\frac a2\right)^2+\left(y+\frac b2\right)^2+\left(z+\frac c2\right)^2 = \frac{a^2+b^2+c^2}{4}-d
\]
a un second membre \textbf{strictement positif} (sinon : point unique si nul, ensemble vide si négatif).
\end{propriete}

\begin{exemple}[Reconnaître et déterminer centre/rayon]
Soit $(S):x^2+y^2+z^2-4x+2y-6=0$. Déterminons si c'est une sphère.

\smallskip
Mise sous forme canonique : $(x-2)^2-4+(y+1)^2-1+z^2-6=0 \iff (x-2)^2+(y+1)^2+z^2=11$.

\smallskip
\textbf{Conclusion} : $(S)$ est la sphère de centre $\Omega(2,-1,0)$ et de rayon $R=\sqrt{11}$.
\end{exemple}

\begin{propriete}[Intersection sphère-plan]
Soit $(S)$ sphère de centre $\Omega$, rayon $R$, et $\mathcal{P}$ un plan. Posons $d=d(\Omega,\mathcal{P})$.
\begin{itemize}[leftmargin=6mm, itemsep=0.5mm]
  \item Si $d>R$ : $(S)\cap\mathcal{P}=\varnothing$.
  \item Si $d=R$ : $\mathcal{P}$ est \textbf{tangent} à $(S)$ en un unique point (le projeté orthogonal de $\Omega$ sur $\mathcal{P}$).
  \item Si $d<R$ : $(S)\cap\mathcal{P}$ est un \textbf{cercle} de centre le projeté orthogonal de $\Omega$ sur $\mathcal{P}$, de rayon $r=\sqrt{R^2-d^2}$.
\end{itemize}
\end{propriete}

\begin{demo}
Soit $H$ le projeté orthogonal de $\Omega$ sur $\mathcal{P}$, $d=\Omega H$. Pour $M\in\mathcal{P}$ : $\Omega M^2=\Omega H^2+HM^2$ (Pythagore, car $(\Omega H)\perp\mathcal{P}$ donc $\perp(HM)$). $M\in(S)\cap\mathcal{P} \iff \Omega M=R \iff HM^2=R^2-d^2$. Si $R^2-d^2>0$, ceci décrit un cercle de centre $H$, rayon $\sqrt{R^2-d^2}$ (dans le plan $\mathcal{P}$).
\end{demo}

\begin{exemple}[Étude d'une intersection sphère-plan]
Soit $(S)$ de centre $\Omega(0,0,0)$, rayon $R=5$, et $\mathcal{P}:x+2y+2z-6=0$. Étudions $(S)\cap\mathcal{P}$.

\smallskip
\[
  d(\Omega,\mathcal{P}) = \frac{|0+0+0-6|}{\sqrt{1+4+4}} = \frac{6}{3} = 2.
\]
Comme $d=2<R=5$ : $(S)\cap\mathcal{P}$ est un \textbf{cercle} de rayon $r=\sqrt{25-4}=\sqrt{21}$, centré au projeté orthogonal de $\Omega$ sur $\mathcal{P}$.
\end{exemple}

\begin{propriete}[Plan tangent à une sphère]
Le plan tangent à la sphère $(S)$ de centre $\Omega$ en un point $A\in(S)$ est le plan passant par $A$, de vecteur normal $\overrightarrow{\Omega A}$.
\end{propriete}

\begin{exemple}[Équation du plan tangent]
Soit $(S)$ de centre $\Omega(1,1,1)$, rayon $3$, et $A(4,1,1)\in(S)$ (vérification : $\Omega A=\sqrt{9+0+0}=3$ $\checkmark$). Déterminons le plan tangent en $A$.

\smallskip
$\overrightarrow{\Omega A}(3,0,0)$, normal au plan tangent. Équation : $3(x-4)+0(y-1)+0(z-1)=0\iff x=4$.
\end{exemple}

% ------------------- SECTION 6 -------------------
\section{Exercices corrigés -- niveau examen national SM}

\begin{exemple}[Exercice 1 -- droite définie par deux points, position par rapport à un plan]
Soient $A(1,2,0)$, $B(3,0,2)$, et $\mathcal{P}:x-y+z-1=0$. Déterminer la position relative de $(AB)$ et $\mathcal{P}$.

\smallskip
\textbf{Corrigé.} $\vec u=\overrightarrow{AB}(2,-2,2)$, normal $\vec n(1,-1,1)$ de $\mathcal{P}$.
\[
  \vec u\cdot\vec n = 2\times1+(-2)\times(-1)+2\times1 = 2+2+2 = 6 \ne 0.
\]
\textbf{Conclusion} : $(AB)$ est \textbf{sécante} à $\mathcal{P}$ (pas parallèle).
\end{exemple}

\begin{exemple}[Exercice 2 -- sphère passant par des points donnés]
Déterminer l'équation de la sphère de centre $\Omega(1,-2,3)$ passant par $A(4,-2,3)$.

\smallskip
\textbf{Corrigé.} $R=\Omega A=\sqrt{(4-1)^2+0+0}=\sqrt9=3$.
\[
  (S):(x-1)^2+(y+2)^2+(z-3)^2=9.
\]
\end{exemple}

\begin{exemple}[Exercice 3 -- intersection droite-sphère]
Soit $(S):x^2+y^2+z^2=9$ et $\mathcal{D}:\begin{cases}x=t\\y=t\\z=1\end{cases}$. Déterminer $(S)\cap\mathcal{D}$.

\smallskip
\textbf{Corrigé.} Substituons : $t^2+t^2+1=9\iff2t^2=8\iff t^2=4\iff t=\pm2$.

\smallskip
\textbf{Conclusion} : deux points d'intersection, $M_1(2,2,1)$ et $M_2(-2,-2,1)$.
\end{exemple}

\begin{exemple}[Exercice 4 -- problème combiné plans et sphère]
Soit $(S):x^2+y^2+z^2-2x+4y-4=0$ et $\mathcal{P}:x-2y+2z+5=0$. Le plan $\mathcal{P}$ coupe-t-il la sphère $(S)$ ? Si oui, préciser le rayon du cercle d'intersection.

\smallskip
\textbf{Corrigé.} Forme canonique de $(S)$ : $(x-1)^2-1+(y+2)^2-4+z^2-4=0\iff(x-1)^2+(y+2)^2+z^2=9$. Centre $\Omega(1,-2,0)$, rayon $R=3$.
\[
  d(\Omega,\mathcal{P}) = \frac{|1-2(-2)+2(0)+5|}{\sqrt{1+4+4}} = \frac{|1+4+5|}{3} = \frac{10}{3}.
\]
Comme $\frac{10}3\approx3{,}33>R=3$ : \textbf{le plan ne coupe pas la sphère}, $(S)\cap\mathcal{P}=\varnothing$.
\end{exemple}

% ------------------- SECTION 7 -------------------
\section{Méthodes et pièges : la boîte à outils de l'examen}

\subsection{Ce qu'on te demande, ce que tu fais}

\begin{center}
\renewcommand{\arraystretch}{1.45}
\sffamily\small
\begin{tabular}{>{\raggedright\arraybackslash}p{7.2cm} >{\raggedright\arraybackslash}p{8cm}}
\rowcolor{qtealdark} \color{white}\bfseries Ce qu'on te demande & \color{white}\bfseries Ton réflexe \\
\rowcolor{tealpale} Position relative de deux droites & Colinéarité des vecteurs directeurs ; sinon, coplanarité via produit mixte. \\
Position relative droite-plan & Comparer $\vec u\cdot\vec n$ à $0$ (parallèle si nul). \\
\rowcolor{tealpale} Intersection de deux plans sécants & Vecteur directeur $\vec n\wedge\vec n'$, puis chercher un point particulier. \\
Reconnaître une sphère & Forme canonique, vérifier que le second membre est strictement positif. \\
\rowcolor{tealpale} Étudier $(S)\cap\mathcal{P}$ & Comparer $d(\Omega,\mathcal{P})$ à $R$ (cercle si $d<R$, tangence si $d=R$, vide si $d>R$). \\
Plan tangent en un point $A$ de la sphère & Normal $=\overrightarrow{\Omega A}$, plan passant par $A$. \\
\rowcolor{tealpale} Intersection droite-sphère & Substituer la paramétrisation de la droite dans l'équation de la sphère. \\
Calculer un angle (droites, plans, droite-plan) & Cosinus (ou sinus pour droite-plan) avec valeur absolue, à partir des vecteurs directeurs/normaux. \\
\end{tabular}
\end{center}

\subsection{Les pièges qui coûtent des points}

\begin{remarque}[Erreurs relevées chaque année par les correcteurs]
\begin{itemize}[leftmargin=6mm, itemsep=0.5mm]
  \item Oublier de vérifier la coplanarité (produit mixte) avant de conclure que deux droites sont sécantes.
  \item Confondre les conditions de parallélisme droite-plan ($\vec u\cdot\vec n=0$) et plan-plan ($\vec n,\vec n'$ colinéaires).
  \item Erreur de signe lors de la mise sous forme canonique d'une équation de sphère (division par $2$ mal appliquée).
  \item Oublier de vérifier que le second membre de la forme canonique est \textbf{positif} avant de conclure qu'il s'agit bien d'une sphère.
  \item Confondre le rayon du cercle d'intersection sphère-plan ($\sqrt{R^2-d^2}$) avec le rayon $R$ de la sphère elle-même.
  \item Oublier qu'un vecteur normal à un plan est aussi vecteur \textbf{directeur} de toute droite orthogonale à ce plan (et inversement).
\end{itemize}
\end{remarque}

% ------------------- RESUME -------------------
\vspace{4mm}
\begin{tcolorbox}[
  enhanced, boxrule=0pt, arc=2mm,
  interior style={left color=qtealdark, right color=qteal},
  left=6mm, right=6mm, top=4mm, bottom=4mm,
  fontupper=\color{white}\sffamily,
  title={\sffamily\bfseries\color{white}\ding{72}~~À retenir},
  colbacktitle=qtealdark, coltitle=white, boxrule=0pt
]
\begin{itemize}[leftmargin=6mm, label={\color{qamber}\ding{212}}, itemsep=1mm]
  \item Droite $(A,\vec u)$ : $\begin{cases}x=x_0+ta\\y=y_0+tb\\z=z_0+tc\end{cases}$ ; plan $(A,\vec n)$ : $a(x-x_0)+b(y-y_0)+c(z-z_0)=0$.
  \item Positions relatives : droites (colinéaires/coplanaires) ; plans (normales colinéaires) ; droite-plan ($\vec u\cdot\vec n$).
  \item Distance point-droite : $\frac{\|\overrightarrow{AM_0}\wedge\vec u\|}{\|\vec u\|}$ ; point-plan : $\frac{|ax_0+by_0+cz_0+d|}{\sqrt{a^2+b^2+c^2}}$.
  \item Sphère $(\Omega,R)$ : $(x-x_0)^2+(y-y_0)^2+(z-z_0)^2=R^2$ ; forme canonique pour reconnaître.
  \item $(S)\cap\mathcal{P}$ selon $d(\Omega,\mathcal{P})$ vs $R$ : vide, tangent, ou cercle de rayon $\sqrt{R^2-d^2}$.
  \item Angles : $\cos\theta=\frac{|\vec u\cdot\vec u'|}{\|\vec u\|\|\vec u'\|}$ (droites/plans, via directeurs/normaux) ; $\sin\theta=\frac{|\vec u\cdot\vec n|}{\|\vec u\|\|\vec n\|}$ (droite-plan).
\end{itemize}
\end{tcolorbox}

\end{document}
